\chapter*{Abstract}
\addcontentsline{toc}{chapter}{Abstract}
\markboth{Abstract}{Abstract}

This End-of-Studies Project (PFE) was carried out in collaboration with
\textbf{BVIRAL}, a company specialized in viral video licensing. The initial
goal was practical: replace disconnected internal tools with one
production-ready platform that supports the full rights workflow.

The project delivered the \textbf{BVIRAL Platform}, built around two
complementary products: (1) a video search portal for discovering clips through
keywords, filters, and natural-language queries, and (2) an Authentic Rights
Portal for onboarding, subscription management, and licensing operations for
creators and media partners.

The platform follows a \textbf{microservices monorepo} architecture with five
independent services: AR-API (FastAPI, PostgreSQL, Celery, Stripe, Clerk),
Videos Search API (Flask, OpenSearch), Agentic Helper API (FastAPI, LangChain,
Qdrant, LangSmith), Video Search Portal (Next.js 15), and Authentic Rights
Portal (Next.js 16). Services are containerized and orchestrated with Docker
Compose.

A core contribution is the \textbf{Agentic Helper}, which combines RAG over
\textit{Qdrant}, LangSmith-managed prompts and tracing, a ReAct video-search
agent, a two-level routing classifier, and Whisper-based audio
transcription. This AI layer improves video discovery, provides contextual
support directly inside the user interfaces, and exposes admin endpoints for
managing FAQ vectors.

\bigskip
\noindent\textbf{Keywords:} Microservices; Next.js; FastAPI; Flask;
OpenSearch; LangChain; RAG; Qdrant; LangSmith; OpenAI; Stripe;
Clerk; Docker; Celery; Video search; Rights management.

\newpage
\chapter*{Summary}
\addcontentsline{toc}{chapter}{Summary}
\markboth{Summary}{Summary}

This report presents the end-to-end design and implementation of the BVIRAL
Platform as a unified system for content rights operations and AI-assisted
video discovery. The architecture separates concerns into independent
microservices, with dedicated APIs for onboarding, payments, search, and
conversational assistance.

From a software engineering perspective, the work focused on production
constraints: clear service boundaries, asynchronous background processing,
containerized deployment, and reproducible CI/CD. From an AI perspective, the
work combined retrieval-based grounding, tool-using agents, intent routing, and
speech transcription to support practical user interactions rather than demo
flows.

The resulting platform is functional in production scenarios and provides a
clear baseline for future upgrades such as multimodal retrieval, stronger LLM
evaluation, and expanded operational tooling.
